\chapter{Distributed Time Synchronization}

To guarantee chronological integrity across the bitsmasher.net distributed Minecraft cluster, system times must be synchronized with sub-millisecond precision. Accurate time tracking prevents write-collisions inside the neural storage databases and ensures exact sequence-locking of spatial telemetry events.

\section{NTP Architecture \& Topology}
The cluster implements a secure, hardened NTP environment to distribute stable clock signals across all execution layers:

\begin{description}
    \item[Authority Node] \texttt{time.lab.bitsmasher.net} (Stratum 1 Local Time Source).
    \item[Daemon Implementation] \texttt{ntpsec} (NTP Security Project).
    \item[Synchronization Mode] Broadcast Client. Subordinate cluster nodes operate in broadcast client mode, listening passively for synchronization packets broadcasted on the local subnet to minimize polling-related network overhead.
    \item[Target Subnet Network] \texttt{10.10.8.0/21}
\end{description}

\section{Operational Security \& Access Control}
To mitigate potential NTP amplification attack vectors and prevent unauthorized time manipulation, strict access control lists are configured within the \texttt{ntpsec} context:

\begin{itemize}
    \item \textbf{Global Restrict Rules:} Default IPv4 traffic is restricted to basic time exchange with zero query privileges:
    \begin{lstlisting}[style=bashstyle]
restrict -4 default kod nomodify notrap nopeer noquery
    \end{lstlisting}
    \item \textbf{Cluster Interrogation Access:} Nodes within the internal \texttt{10.10.8.0/21} range are granted interrogation privileges, allowing the Stargate Control Plane to dynamically audit peer clock health.
    \item \textbf{High-Speed Handshakes:} Authentication is explicitly disabled (\texttt{disable auth}) for local broadcast client operations to ensure high-speed, zero-latency handshakes within the trusted lab environment.
\end{itemize}

\section{System Auditing \& Metrics Loggers}
Performance statistics are compiled daily inside the \texttt{/var/log/ntpstats/} namespace. System administrators must track three primary operational telemetry matrices:
\begin{enumerate}
    \item \textbf{loopstats:} Tracks the local system clock stability, skew, and frequency drift over time.
    \item \textbf{peerstats:} Logs upstream host metrics, capturing jitter, dispersion, and absolute round-trip delay.
    \item \textbf{clockstats:} Records low-level hardware clock subsystem performance metrics.
\end{enumerate}

\section{Actionable Verification Commands}
Synchronization states can be queried interactively from any node (such as Skynet, Edge-T, or the Jetson array) using standard control queries:

\begin{lstlisting}[style=bashstyle, caption={NTP Query Utility Verification}]
# Check peer synchronization status, offset, and jitter
ntpq -p
\end{lstlisting}

\subsection{Thermal Drift Analysis}
Frequency correction offsets are continuously appended to:
\begin{center}
    \texttt{/var/lib/ntp/ntp.drift}
\end{center}

Abrupt changes or high rates of change within this drift file are operational indicators of thermal fluctuations on local edge hardware, such as the Raspberry Pi 5 (Skynet) or Asus Tinker Edge T (Edge-T), directly impacting the physical crystal oscillators.

\section{Impact on the Emerald Mirror Voxel Pipeline}
Establishing a robust \texttt{ntpsec} broadcast timeline directly hardens structural deployment sequences:
\begin{equation}
    \text{Edge-T Vision Sync} \Longleftrightarrow \text{Stargate RCON Dispatch} \Longleftrightarrow \text{Chonk Block Modification}
\end{equation}

By lock-stepping the absolute system time of all nodes, transaction logs written to the Neural Data Vault (\texttt{registry.parquet}) use reliable, monotonic epoch millisecond timestamps. When a player builds in the ``AI Testing Field,'' the vision classification from Edge-T and the RCON injection from Stargate will be sequence-locked, preventing the ``ghost building'' errors common in high-latency distributed clusters.